\documentclass[11pt,a4paper]{article}

\usepackage[utf8]{inputenc}
\usepackage[T1]{fontenc}
\usepackage[scaled=0.92]{helvet}
\usepackage{mathpazo}
\usepackage[margin=1in]{geometry}
\usepackage{amsmath,amssymb}
\usepackage{enumitem}
\usepackage{titlesec}
\usepackage{fancyhdr}
\usepackage{xcolor}
\usepackage{listings}
\usepackage{tabularx}
\usepackage[hidelinks]{hyperref}

% --- SETUP ---
\setlength{\parindent}{0pt}
\setlength{\parskip}{0.5em}
\setlist[enumerate]{leftmargin=*, label=\textbf{\arabic*.}, itemsep=1em}
\setlist[itemize]{leftmargin=*, itemsep=0.2em}

% --- STYLES ---
\definecolor{primary}{HTML}{1F3A5F}
\definecolor{correct}{HTML}{1E8449} % Warna hijau untuk jawaban
\definecolor{codebg}{HTML}{F4F5F7}

\titleformat{\section}{\normalfont\large\bfseries\sffamily\color{primary}}{}{0em}{}[\vspace{-0.5em}\rule{\linewidth}{0.5pt}]
\titleformat{\subsection}{\normalfont\normalsize\bfseries\sffamily\color{primary}}{}{0em}{}

\lstset{
    basicstyle=\ttfamily\small,
    keywordstyle=\color{primary}\bfseries,
    backgroundcolor=\color{codebg},
    frame=none,
    breaklines=true,
    aboveskip=4pt,
    belowskip=4pt,
    xleftmargin=1em,
    columns=fullflexible,
    keepspaces=true
}

\pagestyle{fancy}
\fancyhf{}
\fancyhead[L]{\small\sffamily\bfseries KUNCI JAWABAN EVALUASI KOMPREHENSIF (P1--P7)}
\fancyhead[R]{\small\sffamily TRPL - Politeknik Negeri Batam}
\renewcommand{\headrulewidth}{0.4pt}
\renewcommand{\headrule}{\hbox to\headwidth{\color{gray!50}\leaders\hrule height \headrulewidth\hfill}}

\begin{document}

\begin{center}
	\Large\textbf{\sffamily KUNCI JAWABAN \& RUBRIK PENILAIAN} \\
	\large\textbf{\sffamily UJIAN EVALUASI KOMPREHENSIF BASIS DATA}
\end{center}

\vspace{1em}

\section*{Bagian A: Pilihan Ganda Analitis (Bobot: 30\% / 5 Poin per Soal)}

\begin{enumerate}
	\item \textbf{Jawaban: C} (Total Participation dan Disjoint) \\
	      \textit{Penjelasan:} Total Participation (garis ganda) karena setiap Sivitas \textbf{wajib} terdaftar sebagai Dosen atau Mahasiswa. Disjoint (simbol 'd' atau garis pisah) karena aturan bisnis melarang seseorang menjadi Dosen dan Mahasiswa sekaligus.

	\item \textbf{Jawaban: B} (Karena Keterlibatan Proyek adalah entitas mandiri...) \\
	      \textit{Penjelasan:} Agregasi (P3) digunakan ketika sebuah relasi (Bekerja\_Pada) memiliki relasi lain dengan entitas berbeda (Peralatan). Relasi Bekerja\_Pada diperlakukan sebagai entitas utuh agar dapat dihubungkan dengan Peralatan.

	\item \textbf{Jawaban: C} (Membuat tabel baru yang berisi PK dari A dan B) \\
	      \textit{Penjelasan:} Relasi N:M tidak bisa diselesaikan hanya dengan menaruh FK di salah satu tabel. Mutlak harus membuat tabel *Bridge* (penghubung) yang menampung PK kedua entitas sebagai *Composite Primary Key* (P4/P5).

	\item \textbf{Jawaban: C} (\texttt{NUMERIC(9,2)}) \\
	      \textit{Penjelasan:} Tipe data uang/desimal (P6) tidak boleh menggunakan \texttt{FLOAT} atau \texttt{DOUBLE} karena rentan terhadap *floating-point error* (pembulatan). \texttt{NUMERIC} atau \texttt{DECIMAL} menjamin presisi angka secara eksak.

	\item \textbf{Jawaban: B} (7) \\
	      \textit{Penjelasan:} Fungsi agregat \texttt{COUNT(Nama\_Kolom)} akan \textbf{mengabaikan} baris yang bernilai \texttt{NULL}. Jika ada 10 baris dan 3 baris bernilai \texttt{NULL}, maka hasilnya \(10 - 3 = 7\). (Berbeda dengan \texttt{COUNT(*)} yang akan menghasilkan 10).

	\item \textbf{Jawaban: (Tidak ada yang tepat di opsi, namun secara logika menguji pemahaman \texttt{\_})} \\
	      \textit{Penjelasan:} Wildcard \texttt{\_} berarti \textbf{tepat satu karakter apapun}. Query \texttt{LIKE '\_a\%'} berarti mencari string dengan 2 huruf pertama adalah "karakter apapun" diikuti huruf 'a'.
	      \begin{itemize}
		      \item Al Pacino (A-l) $\rightarrow$ Salah
		      \item Sean Penn (S-e) $\rightarrow$ Salah
		      \item Leonardo (L-e) $\rightarrow$ Salah
		      \item Tom Hanks (T-o) $\rightarrow$ Salah
	      \end{itemize}
	      \textit{Rubrik:} Berikan poin penuh (5) jika mahasiswa menganalisis bahwa semua opsi salah namun menjelaskan logika \texttt{\_}. Berikan poin 2 jika mereka memilih acak tanpa alasan yang jelas.
\end{enumerate}

\vspace{1em}
\section*{Bagian B: Esai \& Studi Kasus (Bobot: 70\%)}

\subsection*{7. EERD to Relational Schema (20 Poin)}
\textbf{Strategi Single Table Inheritance (Menggabungkan Semua Entitas):}
Karena menggunakan strategi Single Table, seluruh atribut dari Supertype dan Subtype digabungkan menjadi satu tabel besar. Atribut yang tidak dimiliki oleh suatu baris akan diisi nilai \texttt{NULL}.

\textbf{Struktur Tabel:}
\begin{itemize}
	\item \textbf{Table:} \texttt{Sivitas}
	\item \textbf{Primary Key:} \texttt{ID\_Sivitas}
	\item \textbf{Atribut Supertype:} \texttt{Nama}, \texttt{Email}
	\item \textbf{Atribut Dosen:} \texttt{NIDN}, \texttt{Gelar}
	\item \textbf{Atribut Mahasiswa:} \texttt{NIM}, \texttt{Tahun\_Masuk}
	\item \textbf{Atribut Mahasiswa S1:} \texttt{SKS}, \texttt{IPK}
	\item \textbf{Atribut Mahasiswa S2:} \texttt{Judul\_Tesis}, \texttt{Pembimbing}
	\item \textbf{Discriminator / Pembeda (Opsional tapi nilai plus):} \texttt{Tipe\_Sivitas} (untuk membedakan apakah dia Dosen, Mhs S1, atau Mhs S2).
\end{itemize}
\textit{Rubrik:}
\begin{itemize}
	\item 5 Poin: Menyebutkan Nama Tabel dan PK dengan benar.
	\item 5 Poin: Menggabungkan atribut Supertype dan Subtype level 1 (Dosen/Mhs).
	\item 10 Poin: Menggabungkan seluruh atribut Subtype (termasuk S1 dan S2) ke dalam satu tabel.
\end{itemize}

\subsection*{8. DDL SQL (20 Poin)}
\begin{lstlisting}[language=SQL]
CREATE TABLE atmwithdrawal (
    transaction_id INT PRIMARY KEY,
    account_number CHAR(3),
    amount NUMERIC(9,2) NOT NULL,
    withdraw_date DATE,
    FOREIGN KEY (account_number)
        REFERENCES account(number)
        ON DELETE CASCADE,
    CHECK (amount >= 0)
);
\end{lstlisting}
\textit{Rubrik:}
\begin{itemize}
	\item 5 Poin: Deklarasi tabel, kolom, dan tipe data tepat.
	\item 5 Poin: \texttt{PRIMARY KEY} pada \texttt{transaction\_id} dan \texttt{NOT NULL} pada amount.
	\item 5 Poin: \texttt{FOREIGN KEY} dan \texttt{ON DELETE CASCADE} ditulis dengan benar.
	\item 5 Poin: Constraint \texttt{CHECK (amount >= 0)} ditulis dengan benar.
\end{itemize}

\subsection*{9. DML SQL (15 Poin)}
\textbf{a. Menampilkan Nama dan Saldo Nasabah Checking di atas 1000 (7.5 Poin)}
\begin{lstlisting}[language=SQL]
SELECT customer.name, account.balance
FROM customer
JOIN account ON customer.id = account.cust_id
WHERE account.type = 'checking'
  AND account.balance > 1000;
\end{lstlisting}
\textit{Rubrik:} 3 Poin \texttt{JOIN ON} yang benar, 2.5 Poin \texttt{WHERE} klausa, 2 Poin \texttt{SELECT} kolom.

\textbf{b. Menampilkan Total Penarikan Bulan April 2008 (7.5 Poin)}
\begin{lstlisting}[language=SQL]
SELECT SUM(amount)
FROM atmwithdrawal
WHERE withdraw_date BETWEEN '2008-04-01' AND '2008-04-30';
-- ATAU menggunakan LIKE:
-- WHERE withdraw_date LIKE '2008-04-%';
\end{lstlisting}
\textit{Rubrik:} 3 Poin fungsi \texttt{SUM()}, 4.5 Poin filter tanggal yang tepat.

\subsection*{10. Analisis Query (15 Poin)}
\textbf{a. Apa hasil logis dari query tersebut? (5 Poin)}\\
Query tersebut akan menampilkan daftar nama aktor (\texttt{MovieStar.name}) yang pernah membintangi suatu film, dengan syarat nama aktor tersebut ada di tabel \texttt{StarsIn}, dan judul filmnya ada di tabel \texttt{Movie}. (Menghubungkan 3 tabel menggunakan \texttt{WHERE} sebagai jembatan relasi).

\textbf{b. Apa yang terjadi jika klausa WHERE dihapus? (10 Poin)}\\
Jika klausa \texttt{WHERE} dihapus, hasilnya adalah **Cartesian Product** (Perkalian Kartesian / Cross Join).
Sistem akan mengalikan setiap baris di tabel \texttt{MovieStar} dengan setiap baris di tabel \texttt{StarsIn}, lalu dikalikan lagi dengan setiap baris di tabel \texttt{Movie}.
Hasilnya adalah data sampah (tidak berelasi) dalam jumlah yang sangat besar (Misal: 10 Star $\times$ 20 StarsIn $\times$ 15 Movie = 3.000 baris data acak).

\end{document}
